\section{Status and Scope}
\label{sec:status}

\subsection{What is built}

\begin{itemize}
  \item The full pipeline: detection, registry dispatch, the
        \texttt{DiagramModule} contract, the \texttt{Scene} contract, and the SVG
        renderer (PNG via \texttt{resvg}).
  \item The Mermaid-compatible kinds and the four Triton-native families (tree,
        struct/memory, topology, poster composition) --- roughly three dozen
        registered \texttt{DiagramKind}s in total.
  \item Poster composition with cell spanning and cross-diagram links.
  \item A dozen theme presets over one unified \texttt{ResolvedTheme}.
  \item A test suite of 318 passing tests covering grammars, layout invariants
        (e.g.\ AVL balance, B-tree key bounds, heap order), rendering, and a smoke
        test that renders every \texttt{examples/**/*.mmd} to valid SVG.
\end{itemize}

\subsection{What is explicitly out of scope}

These are \emph{not} implemented and are not described elsewhere in this document
as if they were:

\begin{itemize}
  \item \textbf{No natural-language input.} Triton compiles source; it does not
        construct diagrams from prompts.
  \item \textbf{No data ingestion.} There are no source adapters that crawl
        external systems into an IR.
  \item \textbf{No PPTX / PDF / HTML / Skia backends.} The single output is SVG
        (rasterised to PNG when required).
  \item \textbf{No general animation pipeline.} The only animated effects are two
        SVG path overlays --- marching ants (\texttt{march}) and particles
        (\texttt{particle}).
  \item \textbf{No agent/MCP ingestion API} and \textbf{no multi-package
        monorepo} --- the project is one package.
\end{itemize}

\subsection{Repository layout}

\begin{center}
\begin{tabular}{ll}
\toprule
Path & Contents \\
\midrule
\texttt{src/contracts/} & Scene, DiagramModule, anchors/ports, theme contracts \\
\texttt{src/frontend/}  & detect + registry dispatch \\
\texttt{src/diagrams/}  & one folder per kind (grammar, IR, layout); the
                          data-structure families are grouped under
                          \texttt{src/diagrams/ds/} (\texttt{struct}, \texttt{tree},
                          \texttt{queue}) \\
\texttt{src/graph/}     & layered, tree, connect layout kernels \\
\texttt{src/scene/}     & Pen builder, cell-strip builder \\
\texttt{src/style/}     & cost-tier styling \\
\texttt{src/theme/}     & presets + resolver \\
\texttt{src/render/}    & the SVG renderer \\
\texttt{examples/}      & \texttt{.mmd} sources beside rendered \texttt{.svg} \\
\bottomrule
\end{tabular}
\end{center}

\subsection{Possible future work}

Honest candidates, none promised: a vector PDF backend, richer auto-layout for
dense graphs, and a first-class JSON/YAML authoring schema to make the structured
input path as ergonomic as the text one.
